\section{Panoramica dell'architettura}\label{sec:panoramica}

Il sistema si articola in due processi distinti: un processo di \emph{indicizzazione}, eseguito in anticipo e
ripetuto solo quando i documenti cambiano, che trasforma la documentazione dell'ateneo in una base di conoscenza
interrogabile; e un processo di \emph{risposta}, eseguito a ogni domanda dello studente, che interroga quella base e
costruisce la risposta. I due processi condividono un unico punto di contatto, l'archivio in cui l'indicizzazione
deposita i frammenti di testo e da cui la risposta li recupera: è questa separazione a rendere la conoscenza del
sistema aggiornabile senza intervenire sul modello linguistico, secondo il principio del RAG discusso nella
Sezione~\ref{sec:rag}. A questi due processi se ne affianca un terzo, la \emph{valutazione}, che misura la qualità
del sistema su un insieme di domande di test.

\subsection*{Componenti}

Il sistema è composto da sette componenti, rappresentati nella Figura~\ref{fig:componenti} e corrispondenti ai
moduli in cui è organizzato il codice.

\begin{figure}[H]
  \centering
  \begin{tikzpicture}[
      comp/.style={draw, thick, rounded corners=3pt, align=center, minimum width=2.6cm, minimum height=0.95cm,
                   font=\footnotesize\bfseries, fill=blue!7},
      lab/.style={font=\scriptsize, align=left, text width=2.6cm, text=black!75},
      labL/.style={lab, align=right},
      arr/.style={-{Stealth[length=2mm]}, thick, black!70},
      head/.style={font=\small\bfseries},
      bubble/.style={draw, rectangle callout, rounded corners=2pt, fill=yellow!12, font=\scriptsize\itshape,
                     align=left, inner sep=3pt},
      card/.style={draw, rounded corners=2pt, fill=white, font=\scriptsize, align=left, inner sep=3pt},
      % icona: pagina web
      webpage/.pic={
        \draw[rounded corners=1pt, fill=white] (-0.45,-0.35) rectangle (0.45,0.35);
        \fill[blue!30, rounded corners=1pt] (-0.45,0.22) rectangle (0.45,0.35);
        \foreach \x in {-0.38,-0.31,-0.24} \fill[white] (\x,0.285) circle (0.022);
        \fill[blue!15] (-0.37,-0.05) rectangle (-0.05,0.14);
        \foreach \y in {0.11,0.03,-0.05} \draw[gray!60, line width=0.7pt] (0.02,\y) -- (0.37,\y);
        \foreach \y in {-0.16,-0.25} \draw[gray!60, line width=0.7pt] (-0.37,\y) -- (0.37,\y);
      },
      % icona: documento con angolo piegato e un'etichetta
      doc/.pic={
        \draw[fill=white] (-0.28,-0.37) -- (-0.28,0.37) -- (0.12,0.37) -- (0.28,0.21) -- (0.28,-0.37) -- cycle;
        \draw (0.12,0.37) -- (0.12,0.21) -- (0.28,0.21);
        \foreach \y in {0.22,0.12} \draw[gray!60, line width=0.6pt] (-0.19,\y) -- (0.05,\y);
        \node[fill=#1, text=white, font=\tiny\bfseries, inner sep=1.2pt] at (0,-0.12) {\tikzpictext};
      },
      % icona: frammenti con metadati
      chunks/.pic={
        \foreach \y/\c in {0.26/orange!30, 0/green!25, -0.26/blue!20} {
          \draw[fill=\c, rounded corners=1pt] (-0.45,\y-0.1) rectangle (0.3,\y+0.1);
          \draw[gray!70, line width=0.5pt] (-0.38,\y) -- (0.2,\y);
          \fill[black!55] (0.4,\y) circle (0.05);
        }
      },
      % icona: studente
      student/.pic={
        \fill[black!55] (0,0.2) circle (0.14);
        \fill[black!55] (-0.24,-0.33) -- (-0.24,-0.12) arc (180:0:0.24) -- (0.24,-0.33) -- cycle;
      },
      % icona: insieme di test
      checklist/.pic={
        \draw[fill=white, rounded corners=1pt] (-0.3,-0.37) rectangle (0.3,0.37);
        \foreach \y in {0.2,0,-0.2} {
          \draw[green!50!black, line width=0.8pt] (-0.22,\y) -- (-0.17,\y-0.05) -- (-0.08,\y+0.06);
          \draw[gray!60, line width=0.6pt] (-0.02,\y) -- (0.22,\y);
        }
      },
      % icona: metriche
      bars/.pic={
        \draw[black!60] (-0.35,-0.3) -- (0.35,-0.3);
        \fill[blue!45] (-0.28,-0.3) rectangle (-0.12,0.05);
        \fill[blue!60] (-0.08,-0.3) rectangle (0.08,0.28);
        \fill[blue!35] (0.12,-0.3) rectangle (0.28,-0.05);
      },
    ]

    % ================= Indicizzazione =================
    \node[head] at (0.3,1.35) {Indicizzazione (offline)};
    \pic at (-0.55,0.35) {webpage};
    \pic[pic text=PDF] at (0.55,0.35) {doc=red!70};
    \node[lab, align=center, text width=3.4cm, anchor=north] at (0,-0.15) {pagine del sito dell'ateneo e dei corsi,
      materiali didattici};
    \node[comp] (ing) at (0,-1.75) {Ingestion};
    \node[comp] (chu) at (0,-4.0) {Chunking};
    \node[comp] (emb) at (0,-6.25) {Embedding};
    \node[cylinder, shape border rotate=90, draw, thick, aspect=0.25, minimum height=1.25cm,
          minimum width=1.5cm, fill=gray!15, font=\scriptsize\bfseries] (qd) at (0.3,-8.85) {Database vettoriale};
    \pic[pic text={\ }] at (-1.25,-8.75) {doc=white};
    \node[lab, align=center, text width=1.4cm] at (-1.25,-9.35) {docstore};

    \draw[arr] (0,-0.95) -- (ing);
    \draw[arr] (ing) -- (chu);
    \pic[pic text={\{ \}}] at (0.75,-2.88) {doc=blue!60};
    \node[lab, right] at (1.05,-2.88) {documenti strutturati (JSON): titoli, sezioni, tabelle};
    \draw[arr] (chu) -- (emb);
    \pic at (0.75,-5.13) {chunks};
    \node[lab, right] at (1.25,-5.13) {frammenti con metadati (corso, ambito, URL)};
    \draw[arr] (emb.south) -- (emb.south |- qd.north);
    \node[font=\ttfamily\scriptsize, fill=white, draw=black!25, inner sep=2pt] at (1.55,-7.25)
      {[0.12, -0.37, \dots]};
    \node[lab, right] at (0.15,-7.7) {vettori densi e sparsi};
    \draw[arr] (chu.west) -| (-1.85,-6.0) |- (-1.25,-7.95) -- (-1.25,-8.33);

    % ================= Risposta =================
    \node[head] at (9.4,1.35) {Risposta a una domanda (online)};
    \pic at (9.4,0.35) {student};
    \node[bubble, callout absolute pointer={(9.1,0.4)}] at (6.9,0.55) {le lezioni sono\\obbligatorie?};
    \node[comp] (ui) at (9.4,-1.75) {UI};
    \draw[arr, {Stealth[length=2mm]}-{Stealth[length=2mm]}] (9.4,-0.05) -- (ui);
    \node[comp] (ret) at (9.4,-8.85) {Retrieving};
    \node[comp] (gen) at (9.4,-11.2) {Generation};

    \draw[arr] (ui) -- (ret);
    \node[card, anchor=east] at (9.15,-4.6) {\textbf{domanda:} \emph{le lezioni sono obbligatorie?}\\
      \textbf{corso:} Infermieristica\\ \textbf{cronologia:} turni precedenti};

    \draw[arr] (qd.east) -- node[above, font=\scriptsize] {frammenti candidati} (ret.west);

    \draw[arr] (ret) -- (gen);
    \node[card, anchor=east, text width=4.1cm] at (9.15,-10.03) {\textbf{[1]} \emph{Frequenza delle lezioni}:
      \emph{``La frequenza ai corsi è obbligatoria per il 66\% delle ore\dots''}};

    \draw[arr] (gen.east) -- ++(0.55,0) |- (ui.east);
    \node[bubble, callout absolute pointer={(10.72,-5.4)}, anchor=west, text width=2.45cm] at (11.05,-5.4)
      {Sì, la frequenza è obbligatoria per il 66\% delle ore di lezione.\\[2pt]
       \upshape\textbf{Fonti:} Frequenza delle lezioni};

    % ================= vLLM =================
    \node[draw, thick, rounded corners=3pt, fill=violet!12, align=center, minimum width=2.0cm,
          minimum height=1.3cm, font=\footnotesize] (vllm) at (12.45,-9.4)
      {\textbf{vLLM}\\[1pt]{\scriptsize modello linguistico}\\[-1pt]{\scriptsize su GPU}};
    \foreach \x/\t in {11.85/1, 12.45/2, 13.05/3} {
      \node[draw, rounded corners=1pt, fill=white, font=\tiny, minimum width=0.42cm, minimum height=0.32cm,
            inner sep=1pt] (req\t) at (\x,-8.05) {\t};
      \draw[-{Stealth[length=1.5mm]}, violet!70!black] (req\t.south) -- (\x,-8.72);
    }
    \node[font=\scriptsize, text=violet!60!black, align=center, text width=1.9cm] at (12.5,-7.45)
      {richieste concorrenti};
    \draw[{Stealth[length=2mm]}-{Stealth[length=2mm]}, thick, violet!70!black] (ret.east) -- (ret.east -| vllm.west);
    \draw[{Stealth[length=2mm]}-{Stealth[length=2mm]}, thick, violet!70!black]
      (gen.south) -- ++(0,-0.75) -| (vllm.south);

    % ================= Valutazione =================
    \node[head] at (0.3,-10.25) {Valutazione};
    \node[comp, fill=gray!8] (eva) at (0,-11.2) {Evaluation};
    \pic at (-0.95,-12.6) {checklist};
    \node[lab, align=center, text width=2.2cm] at (-0.95,-13.55) {domande con risposte e fonti attese};
    \pic at (0.95,-12.6) {bars};
    \node[lab, align=center, text width=2.2cm] at (0.95,-13.3) {metriche (JSON)};
    \draw[arr] (-0.95,-12.2) -- (-0.95,-11.68);
    \draw[arr] (0.95,-11.68) -- (0.95,-12.2);
    \draw[arr, dashed, {Stealth[length=2mm]}-{Stealth[length=2mm]}] (eva) --
      node[above, font=\scriptsize] {domande di test} node[below, font=\scriptsize] {risposte e frammenti} (gen);
  \end{tikzpicture}
  \caption{Componenti del sistema e dati che si scambiano, seguendo una domanda di esempio. A sinistra
  l'indicizzazione trasforma pagine web e documenti in frammenti e vettori, salvati nel database vettoriale e nel docstore; a
  destra la domanda dello studente attraversa interfaccia, recupero e generazione, fino alla risposta con la fonte
  citata. In basso la valutazione sottopone alla stessa pipeline le domande di test e ne misura la qualità.
  L'Embedding opera anche al momento della domanda, per rappresentarla nello stesso spazio vettoriale dei
  frammenti. In viola le chiamate al modello linguistico, servito da vLLM: le inviano il Retrieving (analisi della
  domanda, HyDE) e la Generation; vLLM elabora insieme le richieste che arrivano contemporaneamente, anche da
  domande di studenti diversi. Per descrivere le immagini, in fase di indicizzazione l'Ingestion usa una seconda
  istanza di vLLM con un modello multimodale, attiva solo durante la conversione.}
  \label{fig:componenti}
\end{figure}

\paragraph{Ingestion.} Acquisisce i documenti e li porta in un formato strutturato comune. Le pagine web dei corsi
di laurea e di ateneo vengono scaricate da un crawler dedicato; queste e gli altri documenti (PDF, presentazioni)
vengono poi convertiti con Docling in una rappresentazione che conserva titoli, sezioni e tabelle, e corretti dalle
imprecisioni ricorrenti della conversione.

\paragraph{Chunking.} Suddivide ogni documento convertito in frammenti che rispettano la struttura delle sezioni e
associa a ciascuno i metadati necessari al recupero, tra cui il corso di laurea a cui appartiene o l'indicazione che
si tratta di una pagina di ateneo. Conserva inoltre il legame tra i frammenti e la sezione da cui provengono, usato
dalla fusione gerarchica in fase di recupero.

\paragraph{Embedding.} Calcola la rappresentazione vettoriale di ogni frammento, densa e sparsa, e la salva nel
database vettoriale insieme ai metadati; lo stesso modello calcola, al momento della domanda, la rappresentazione
della domanda stessa.

\paragraph{Retrieving.} A ogni domanda individua i frammenti pertinenti: analizza la domanda (ne ricostruisce il
senso rispetto alla conversazione, la scompone se contiene più richieste, riconosce il corso di laurea), cerca i
candidati tra le pagine di ateneo e quelle del corso, li riordina in base alla pertinenza e scarta quelli
insufficienti.

\paragraph{Generation.} Costruisce il prompt con la domanda e i frammenti recuperati, interroga il modello
linguistico e restituisce la risposta con le fonti citate. Quando il recupero non produce frammenti pertinenti, il
modello non viene interpellato per rispondere.

\paragraph{UI.} È l'interfaccia di chat con cui lo studente interagisce e il punto in cui il percorso di una domanda
viene coordinato: mantiene la cronologia della conversazione e il corso dello studente, decide quando chiedere
all'utente di indicare il corso e mostra la risposta man mano che viene generata.

\paragraph{Evaluation.} Esegue lo stesso percorso di recupero e generazione su un insieme di domande di test con
risposta nota e ne misura la qualità, separatamente per il recupero e per la risposta finale. Non fa parte del
sistema in esercizio, ma è lo strumento con cui sono state verificate le scelte descritte in questo capitolo
(Capitolo~\ref{cap:validazione}).

\medskip
\noindent La Tabella~\ref{tab:componenti} riassume, per ciascun componente, le tecnologie adottate e il momento in
cui opera;

\begin{table}[H]
  \centering
  \footnotesize
  \begin{tabularx}{\textwidth}{@{} L{2.2cm} >{\raggedright\arraybackslash}X L{3.4cm} @{}}
    \toprule
    \textbf{Componente} & \textbf{Tecnologie} & \textbf{Quando opera} \\
    \midrule
    Ingestion & crawler dedicato, Docling & indicizzazione \\
    Chunking & Docling, LlamaIndex & indicizzazione \\
    Embedding & BGE-M3, Qdrant & indicizzazione e a ogni domanda \\
    Retrieving & Qdrant, bge-reranker-v2-m3, LlamaIndex, modello linguistico & a ogni domanda \\
    Generation & modello linguistico (Qwen2.5-14B-Instruct) servito da vLLM & a ogni domanda \\
    UI & Chainlit & a ogni domanda \\
    Evaluation & RAGAS e metriche di recupero dedicate & solo in valutazione \\
    \bottomrule
  \end{tabularx}
  \caption{Componenti del sistema, tecnologie adottate e momento in cui operano.}
  \label{tab:componenti}
\end{table}
